%%%%%%%%%%%%%%%%%%%%%%%%%%%%% TCC %%%%%%%%%%%%%%%%%%%%%%%%%%%%%%%%
%
% Template para TCC da Universidade Federal da Paraíba
%
% Adaptado para o projeto de Federated Learning com SDN
% Autor: Cândido Leandro de Queiroga Bisneto
% Orientador: Fernando Menezes Matos
%
%%%%%%%%%%%%%%%%%%%%%%%%%%%%%%%%%%%%%%%%%%%%%%%%%%%%%%%%%%%%%%%%%%%

\documentclass{tcc}

\begin{document}
\pagestyle{empty} %retira numeração da página
%Dados do TCC%
\author{Cândido Leandro de Queiroga Bisneto}
\title{Otimização de Modelos de Aprendizado Federado com Redes Definidas por Software}
\newcommand{\subtitulo}{}
\newcommand{\nomedocurso}{Ciência da Computação}
\newcommand{\titulobar}{Ciência da Computação}
\newcommand{\orientador}{Fernando Menezes Matos}
\newcommand{\profa}{Fernando Menezes Matos}
\newcommand{\profb}{[Nome do Professor B]}
\newcommand{\profc}{[Nome do Professor C]}
\newcommand{\insta}{Universidade Federal da Paraíba}
\newcommand{\instb}{Universidade Federal da Paraíba}
\newcommand{\instc}{Universidade Federal da Paraíba}
\newcommand{\coordenador}{[Nome do Coordenador]}
\newcommand{\departamento}{Ciência da Computação}


% CAPA
\begin{center}
\LARGE{\bf \thetitle}\\
\end{center}

\vspace{1em}

\begin{center}
\theauthor
\end{center}

\vfill

\begin{figure}[H]
\centering
\includegraphics[height=3cm]{figures/ci_logo.png}
\hspace{1cm}
\includegraphics[height=3cm]{figures/ufpb_logo.png}
\end{figure}

\begin{center}
CENTRO DE INFORMÁTICA \\
UNIVERSIDADE FEDERAL DA PARAÍBA
\end{center}

\vspace{2in}

\begin{center}
João Pessoa, 2025
\end{center}
\afterpage{\blankpage \addtocounter{page}{1}}

% FOLHA DE ROSTO
\newpage
\begin{center}
\theauthor
\end{center}
\vspace{3in}
\begin{center}
\LARGE{\thetitle}\\
\end{center}

\vspace{2in}

\begin{flushright}
Monografia apresentada ao curso \nomedocurso \\ do Centro de Informática, da Universidade Federal da Paraíba, \\ como requisito para a obtenção do grau de Bacharel em \titulobar
\\
\vspace{0.2in}

Orientador: \orientador

\end{flushright}

\vfill
\begin{center}
Dezembro de 2025
\end{center}

\newpage

% PÁGINA EM BRANCO (para ficha catalográfica no verso)
$ $
\vfill

\newpage

% PÁGINA DE APROVAÇÃO
\begin{figure}[H]
\centering
\includegraphics[height=3cm]{figures/ci_logo.png}
\hspace{1cm}
\includegraphics[height=3cm]{figures/ufpb_logo.png}
\end{figure}

\begin{center}
CENTRO DE INFORMÁTICA \\
UNIVERSIDADE FEDERAL DA PARAÍBA
\end{center}

\vspace{0.05in}

Trabalho de Conclusão de Curso de \nomedocurso \ intitulado \textit{\bf \em \thetitle} de autoria de \theauthor, aprovada pela banca examinadora constituída pelos seguintes professores: \\

\vspace{0.7in}

\hrule
\noindent Prof. Dr. \profa\\
\insta\\

\vspace{0.25in}

\hrule
\noindent Prof. \profb\\
\instb\\

\vspace{0.25in}

\hrule
\noindent Prof. \profc\\
\instc\\

\vspace{0.8in}

\hrule
\noindent Coordenador(a) do Curso de \departamento\\
\coordenador\\
CI/UFPB\\

\vfill

\begin{center}
João Pessoa, Dezembro de 2025
\end{center}

\afterpage{\blankpage \addtocounter{page}{1}}

% DEDICATÓRIA
\newpage
\section*{\centering{DEDICATÓRIA}}
[Texto da dedicatória - opcional]


% AGRADECIMENTOS
\newpage
\section*{\centering{AGRADECIMENTOS}}

[Texto dos agradecimentos]


% RESUMO
\newpage
\section*{\centering{RESUMO}}

[Inserir resumo em português - um único parágrafo informativo contendo: objetivos do trabalho, justificativa e resultados alcançados. Máximo 200 palavras.]

{\bf Palavras-chave:} Aprendizado Federado, Redes Definidas por Software, XGBoost, LightGBM, CatBoost, Flower Framework, Privacidade de Dados.


% ABSTRACT
\newpage
\section*{\centering{ABSTRACT}}

[Insert abstract in English - one informative paragraph containing: work objectives, justification, and results achieved. Maximum 200 words.]

{\bf Keywords:} Federated Learning, Software-Defined Networking, XGBoost, LightGBM, CatBoost, Flower Framework, Data Privacy.


% LISTA DE FIGURAS
\newpage
\renewcommand{\listfigurename}{\centering LISTA DE FIGURAS}
\listoffigures

% LISTA DE TABELAS
\newpage
\renewcommand{\listtablename}{\centering LISTA DE TABELAS}
\listoftables

% LISTA DE ABREVIATURAS
\newpage
\section*{\centering{LISTA DE ABREVIATURAS}}

API - Application Programming Interface

AUC - Area Under the Curve

CART - Classification and Regression Trees

CPU - Central Processing Unit

DP - Differential Privacy

FL - Federated Learning

GBDT - Gradient Boosting Decision Trees

GPU - Graphics Processing Unit

gRPC - Google Remote Procedure Call

HTTP - Hypertext Transfer Protocol

IID - Independent and Identically Distributed

LGPD - Lei Geral de Proteção de Dados

ML - Machine Learning

non-IID - Non-Independent and Identically Distributed

QoS - Quality of Service

ROC - Receiver Operating Characteristic

SDN - Software-Defined Networking

TCC - Trabalho de Conclusão de Curso

TCP - Transmission Control Protocol

UFPB - Universidade Federal da Paraíba

UDP - User Datagram Protocol


% SUMÁRIO
\newpage
\pagestyle{plain} %mostra numeração da página%
\tableofcontents

% CAPÍTULOS
\newpage
\section{INTRODUÇÃO}

\subsection{Motivação}

\subsection{Objetivo}

\subsection{Estrutura da monografia}
Esta monografia está dividida da seguinte maneira: na segunda seção é feita a revisão da literatura sobre aprendizado federado e redes SDN. Na terceira seção são apresentados os fundamentos teóricos necessários para compreensão do trabalho. Na quarta seção é apresentada a metodologia utilizada no desenvolvimento do trabalho. Na quinta seção o dataset utilizado é descrito em detalhes. Na sexta seção a configuração experimental é apresentada. Na sétima seção os resultados obtidos são apresentados e analisados. Na oitava seção é feita a discussão dos resultados. Na nona seção é feita a conclusão deste trabalho e trabalhos futuros são propostos.

\section{REVISÃO DE LITERATURA}

\subsection{Trabalhos Relacionados sobre Aprendizado Federado}

\subsection{Trabalhos sobre Modelos Baseados em Árvores}

\subsection{Trabalhos sobre SDN e Otimização de Rede}

\subsection{Análise Comparativa dos Trabalhos}

\section{FUNDAMENTOS TEÓRICOS}

\subsection{Aprendizado Federado}

\subsubsection{Definição e Características}

\subsubsection{Estratégias de Agregação}

\subsection{Modelos Baseados em Árvores de Decisão}

\subsubsection{XGBoost}

\subsubsection{LightGBM}

\subsubsection{CatBoost}

\subsection{Redes Definidas por Software (SDN)}

\subsubsection{Arquitetura SDN}

\subsubsection{OpenFlow}

\subsubsection{Aplicações em Aprendizado de Máquina}

\subsection{Framework Flower}

\subsubsection{Arquitetura do Flower}

\subsubsection{Implementação de Clientes e Servidores}

\section{METODOLOGIA}

\subsection{Visão Geral da Solução Proposta}

\subsection{Arquitetura do Sistema}

\subsubsection{Componentes do Sistema}

\subsubsection{Fluxo de Comunicação}

\subsection{Implementação dos Modelos Federados}

\subsubsection{Adaptação de XGBoost para FL}

\subsubsection{Adaptação de LightGBM para FL}

\subsubsection{Adaptação de CatBoost para FL}

\subsection{Integração com SDN}

\subsubsection{Configuração da Rede SDN}

\subsubsection{Políticas de QoS}

\subsubsection{Monitoramento de Tráfego}

\subsection{Tecnologias Utilizadas}

\begin{table}[H]
\centering
\begin{tabularx}{\textwidth}{|l|X|}
\hline
\bf Componente & \bf Tecnologia \\ \hline
Framework FL & Flower 1.6+ \\ \hline
Modelos & XGBoost, LightGBM, CatBoost \\ \hline
Linguagem & Python 3.9+ \\ \hline
SDN Controller & [A definir] \\ \hline
Comunicação & gRPC \\ \hline
\end{tabularx}
\caption{Tecnologias utilizadas no projeto}
\label{tab:tecnologias}
\end{table}

\section{DATASET E PREPARAÇÃO DOS DADOS}

\subsection{Seleção do Dataset}

\begin{table}[H]
\centering
\begin{tabularx}{\textwidth}{|l|X|}
\hline
\bf Característica & \bf Valor \\ \hline
Nome & [Nome do dataset] \\ \hline
Fonte & [URL/Repositório] \\ \hline
Número de amostras & [N] \\ \hline
Número de features & [M] \\ \hline
Classes & [Lista de classes] \\ \hline
Tipo de problema & Classificação binária/multiclasse \\ \hline
\end{tabularx}
\caption{Metadados do dataset}
\label{tab:dataset_metadata}
\end{table}

\subsection{Descrição das Features}

\subsection{Análise Exploratória dos Dados}

\subsection{Pré-processamento}

\subsubsection{Tratamento de Valores Ausentes}

\subsubsection{Normalização e Padronização}

\subsubsection{Codificação de Variáveis Categóricas}

\subsection{Particionamento dos Dados}

\subsubsection{Particionamento IID}

\subsubsection{Particionamento non-IID}

\section{CONFIGURAÇÃO EXPERIMENTAL}

\subsection{Ambiente Computacional}

\begin{table}[H]
\centering
\begin{tabularx}{\textwidth}{|l|X|}
\hline
\bf Componente & \bf Especificação \\ \hline
Processador & [Modelo e frequência] \\ \hline
Memória RAM & [Quantidade] GB \\ \hline
GPU & [Modelo] ou N/A \\ \hline
Sistema Operacional & [Nome e versão] \\ \hline
Python & [Versão] \\ \hline
\end{tabularx}
\caption{Especificações do ambiente computacional}
\label{tab:hardware}
\end{table}

\subsection{Hiperparâmetros dos Modelos}

\subsection{Configuração do Aprendizado Federado}

\subsection{Experimentos Realizados}

\subsection{Métricas de Avaliação}

\section{RESULTADOS}

\subsection{Resultados em Cenário IID}

\subsubsection{Comparação de Modelos}

\subsubsection{Comparação de Estratégias}

\subsection{Resultados em Cenário non-IID}

\subsubsection{Comparação de Modelos}

\subsubsection{Impacto da Heterogeneidade}

\subsection{Análise de Convergência}

\subsection{Impacto da Integração SDN}

\subsection{Análise Estatística}

\subsubsection{Teste de Friedman}

\subsubsection{Teste Post-hoc de Nemenyi}

\section{DISCUSSÃO}

\subsection{Interpretação dos Resultados}

\subsection{Comparação com Estado da Arte}

\subsection{Limitações do Trabalho}

\subsection{Implicações Práticas}

\section{CONCLUSÕES E TRABALHOS FUTUROS}

\subsection{Conclusões}

\subsection{Contribuições}

\subsection{Trabalhos Futuros}


% REFERÊNCIAS
\renewcommand{\refname}{\centering REFERÊNCIAS}
\addcontentsline{toc}{section}{REFERÊNCIAS}
\begin{thebibliography}{99}
% Referências em formato ABNT

% Exemplo de artigo
\bibitem{exemplo2020}AUTOR, A.B.; COAUTOR, C.D. 
Título do artigo. {\bf Nome da Revista}, v.10, n.2, 
p.45-62, 2020.

% Exemplo de livro
\bibitem{exemplo_livro}AUTOR, Nome. 
{\bf Título do Livro}. Editora, Cidade, Ano.

% Exemplo de conferência
\bibitem{exemplo_conf}AUTOR, A.; COAUTOR, B.
Título do artigo. {\bf Nome da Conferência}, 
p.100-110, 2021.

% Adicionar referências conforme necessário

\end{thebibliography}

% ANEXOS
% Anexos (opcional)

% Descomentar se houver anexos

% \section*{\centering{ANEXO A – CÓDIGO FONTE}}

% \subsection*{Implementação do Cliente Flower}
% \begin{lstlisting}[language=Python]
% # Código aqui
% \end{lstlisting}

% \section*{\centering{ANEXO B – DOCUMENTAÇÃO ADICIONAL}}


\end{document}
